\section{End-to-End Harness (Implemented, Not Executed)}
\label{sec:method-e2e-full}

\paragraph{Environment and tasks} Live websites change and block automated traffic, and WebArena-style site images were not available offline, so we built five self-hosted sites served on the loopback interface with fictional content that a model cannot answer from memory: a search engine, an encyclopedia, a shop with search and cart, a town-office contact form and newsletter, and a notes application; a further path redirects to a sign-in wall. The 40 tasks (10 development, 30 test) cover search, navigation, text entry, extraction, form submission, multi-step tasks and 5 unsatisfiable requests, at four difficulty tiers by the number of actions a solution needs (1, 2--3, 4--6, more). The notes site appears only in the test split. Tasks needing an account, payment, a CAPTCHA, the live web, or an outcome that neither server state nor answer text can verify were excluded by design. Because the search-home check $H$ recognizes only public search engines, it never fires on these sites; its ablation in the harness is uninformative.

\paragraph{Oracle} $\Phi^{\star}$ is computed from signals that neither the agent nor $\Phi$ uses: server-side records of the sites (cart contents, form submissions, saved notes, field values autosaved by page scripts), the final URL, and the ground-truth facts in the site data matched against the agent's answer; unsatisfiable tasks are never achieved. It calls no model and shares no code with the verifier. We measured its error rate on 114 scripted trajectories driven directly through Playwright without the agent: a correct solution for each task, an idle trajectory that only opens the start page, and a near miss for 34 tasks (wrong value, wrong item, extra submission, partial query).

\paragraph{Protocol} Table~\ref{tab:e2e} (Section~\ref{app:tasks}) lists the configurations: the full system, two baselines---(a) no completion check and (b) self-verification by the planner model---and one-factor ablations of each conjunct, memory, recovery and each recovery strategy. Each configuration runs in its own process with its own database and memory store, at temperature 0 with a fixed seed, three trials per task, and the same budget of 15 planner calls (after which the recovery engine may grant its remaining attempts). Per episode the harness records the oracle verdict, the agent's status, browser actions, model calls and tokens, wall-clock time, each verification decision with its latency, recovery strategies, retrieved memories, Ollama version and model digests, and the host of every page request Playwright reports. The analysis reports success, FCR and FRR with Wilson intervals, per-category breakdowns, exact McNemar tests and task-level bootstrap intervals~\cite{efron1979bootstrap} for differences to the full system, and the overhead of $\Phi$ relative to baseline (a). A published agent framework as a third baseline would need its own model and prompts and is outside the scope of this paper.


\subsection{Task Set}
\label{app:tasks}

The 40 tasks (\texttt{benchmarks/e2e/tasks.json}) are distributed as follows (development/test): extraction 2/8, form submission 2/7, multi-step 1/5, navigation 2/2, search 1/2, text entry 1/2, unsatisfiable 1/4; tiers 1--4 contain 3/6, 4/11, 3/12 and 0/1 tasks. The test split uses all five sites; the notes site appears only there. The harness (\texttt{run\_e2e.py}) takes the split, the number of trials and a list of the configurations of Table~\ref{tab:e2e}; the analysis script (\texttt{analyze\_e2e.py}) writes result tables with Wilson intervals, McNemar tests and bootstrap intervals.

\begin{table}[t]
\centering
\caption{End-to-end evaluation design (implemented in \texttt{benchmarks/e2e}, not yet run with a model). Each configuration changes one factor of the full system; all use the same tasks, model, sampling and budget.}
\label{tab:e2e}
\footnotesize
\setlength{\tabcolsep}{2.2pt}
\begin{tabular}{@{}p{2.95cm}p{2.85cm}p{2.6cm}@{}}
\toprule
Configuration & Change & Question \\
\midrule
\texttt{full} & -- & task success, FCR, FRR, cost \\
\texttt{no\_verification} & completion accepted wherever $\Phi$ would be evaluated & baseline (a); gain and overhead of $\Phi$ \\
\texttt{self\_verification} & planner model judges its own completion & baseline (b) \\
\texttt{gate\_without\_}$c$ & conjunct $c$ removed & per-evidence ablation \\
\texttt{no\_memory} & memory and RAG off & contribution of memory \\
\texttt{no\_evidence} & no per-action execution records & traceability \\
\texttt{no\_recovery} & engine off; inline retries remain & contribution of recovery \\
\texttt{recovery\_}$s$\texttt{\_only} & strategy $s$ pinned & per-strategy recovery \\
\texttt{single\_model} & routing off & routing \\
\texttt{vision\_only\_\allowbreak grounding} & no element list & grounding \\
\bottomrule
\end{tabular}
\end{table}

